\section{Worker III: recurrence transport with singularity-aware seeds}
\label{sec:recurrences}

\subsection{Why two good recurrence proofs may not meet}

Suppose one worker proves that a sequence $f$ is annihilated by $A$, while a
different derivation proves that $g$ is annihilated by $B$. Matching their first
several values does not automatically establish equality. One needs a common
induction rule and an exact account of the indices where that rule cannot
propagate equality.

Use the forward shift $E$, acting by $(Ef)(n)=f(n+1)$. A polynomial-coefficient
operator is
\[
A=\sum_{i=0}^{r}a_i(n)E^i,
\qquad (Af)(n)=\sum_{i=0}^{r}a_i(n)f(n+i).
\]
The coefficients lie in $\Q[n]$, and all recurrence hypotheses in the prototype
are intended for every natural $n$. The order is the largest index with a
nonzero coefficient polynomial.

\subsection{Noncommutative multiplication is part of the certificate}

Because shifting changes a coefficient, operator multiplication obeys
\begin{equation}\label{eq:ore}
\bigl(a_i(n)E^i\bigr)\bigl(b_j(n)E^j\bigr)
=a_i(n)b_j(n+i)E^{i+j}.
\end{equation}
In particular $En=(n+1)E$, not $nE$. The supplied checker expands this identity
coefficientwise using exact polynomial shift. A dedicated unit test rejects
commutative multiplication as a substitute.

A certificate of a common \emph{left} multiple consists of three operators
$L,U,V$ such that
\begin{equation}\label{eq:left-multiple}
L=UA=VB.
\end{equation}
The problem supplies $A,B$; the certificate cannot redefine them. The checker
requires all three polynomials-in-$E$ to be canonical with nonzero final
coefficients and verifies both products by \eqref{eq:ore}. It does not claim
$L$ is a least common left multiple.

\begin{proposition}[Annihilator transport]
If $Af=0$ and $Bg=0$ on $\N$, and \eqref{eq:left-multiple} holds with polynomial
coefficients and nonnegative shifts, then $Lf=Lg=0$ on $\N$.
\end{proposition}
\begin{proof}
For all $n$, $(UA)f(n)=U(Af)(n)=0$, since every shifted index remains natural.
The analogous calculation applies to $V(Bg)$. Identity \eqref{eq:ore} is exactly
what makes composition of actions agree with operator multiplication.
\end{proof}

The prototype searches for $U,V$ of bounded order and coefficient degree,
forms the linear equations for $UA-VB=0$, and solves them exactly. Any nonzero
solution yields an acceptable $L$. In a production system, \code{ore_algebra}
can supply more sophisticated left-multiple and desingularisation algorithms;
its output should still pass this simpler polynomial identity check~\cite{ore,ore-paper}.
Rational-function operator coefficients would require a denominator policy and
validity-domain ledger. The implemented format excludes them.

\subsection{A concrete common multiple}

Let
\[
A=E-(n+1),\qquad B=E-2(n+1).
\]
The search returns
\[
U=E-2(n+2),\qquad V=E-(n+2),
\]
and
\[
L=E^2-3(n+2)E+2(n+1)(n+2).
\]
Expanding $UA$ and $VB$ with \eqref{eq:ore} verifies the result. Treating $E$
as a commuting indeterminate would produce the wrong middle coefficient.
This example is a general operator-composition test, not a claim that the
nonzero factorial and doubled-factorial solutions of the two first-order
operators are equal. Equality still needs the relevant initial-value proofs.

The certificate also permits a deliberate nonminimal common multiple. This
helps separate discovery from acceptance: if a producer finds a larger $L$
quickly, the checker can accept it without proving an optimality theorem.
The cost is potentially more seed obligations and additional apparent
singularities, which must be accounted for rather than cancelled informally.

\subsection{The singularity problem}

Write a common recurrence as
\begin{equation}\label{eq:common-rec}
\sum_{i=0}^{r}p_i(n)f(n+i)=0,\qquad p_r\not\equiv0,\quad r\ge1.
\end{equation}
At an index $n$ with $p_r(n)\ne0$, the preceding $r$ values determine $f(n+r)$.
At a root of $p_r$, that argument fails. Merely knowing that $p_r$ is a nonzero
polynomial is insufficient: a polynomial can be nonzero as an expression and
zero at a specific natural index.

\begin{example}[A missed singularity invalidates prefix reasoning]
Consider
\[
(n-5)\bigl(f(n+1)-2f(n)\bigr)=0.
\]
The zero sequence and
\[
g(n)=\begin{cases}0,&n<6,\\2^{n-6},&n\ge6\end{cases}
\]
both satisfy it. They agree at indices $0$ through $5$ and disagree at index $6$.
Thus even six matching initial values do not justify equality for this
first-order recurrence unless the singular step is handled. This countermodel
is exercised by the unit tests.
\end{example}

\subsection{A selective seed theorem}

Let
\[
\Sigma=\{s\in\N:p_r(s)=0\},\qquad
I=\{0,1,\ldots,r-1\}\cup\{s+r:s\in\Sigma\}.
\]
Because $p_r$ is a nonzero polynomial over a field, $\Sigma$ is finite.

\begin{theorem}[Singularity-aware recurrence equality]\label{thm:seeds}
Let $f,g:\N\to\Q$ both satisfy \eqref{eq:common-rec}. If $f(i)=g(i)$ for every
$i\in I$, then $f(n)=g(n)$ for every $n\in\N$.
\end{theorem}
\begin{proof}
Set $h=f-g$. It satisfies the same homogeneous recurrence. Use strong induction
on $m$. If $m<r$, then $h(m)=0$ by the seed assumptions. Otherwise write
$m=n+r$. If $p_r(n)=0$, then $m\in I$ and the seed assumption again applies.
If $p_r(n)\ne0$, the recurrence at $n$ reads
\[
p_r(n)h(m)=-\sum_{i=0}^{r-1}p_i(n)h(n+i)=0,
\]
because $n+i<m$ for every $i<r$. Cancellation in the field gives $h(m)=0$.
\end{proof}

This set of seeds is sufficient, not asserted minimal. Additional recurrence
relations can sometimes determine a value that this theorem requests
separately. Its advantage over demanding every value through the largest
singular index is that it asks only for the initial window and the exceptional
pivots. Repeated roots do not cause duplicate seed obligations.

For the example above, the required set is $\{0,6\}$. For a third-order
recurrence with natural leading roots $0,1,5$, it is
$\{0,1,2,3,4,8\}$. A seed at $s$ instead of $s+r$ is an off-by-order error;
it does not repair the missing pivot.

\subsection{A complete finite cover of natural singularities}

The checker must not trust a producer-supplied list of roots. Let
\[
p_r(n)=a_dn^d+\cdots+a_0,\qquad a_d\ne0.
\]
For $d>0$, choose an integer
\begin{equation}\label{eq:cauchy}
B\ge1+\max_{0\le i<d}|a_i/a_d|.
\end{equation}
Every complex root has absolute value at most this bound. For completeness,
if $|z|>1+C$ with $C=\max_{i<d}|a_i/a_d|$, then a zero would imply
\[
|z|^d\le C\sum_{i=0}^{d-1}|z|^i
\le\frac{C}{|z|-1}|z|^d<|z|^d,
\]
a contradiction. The constant-polynomial case has no roots and uses $B=0$.

A singularity-cover certificate supplies $B$, the claimed natural roots, and
the proposed seed set. The checker verifies \eqref{eq:cauchy}, evaluates $p_r$
exactly at every integer in $[0,B]$, compares the complete sorted root list, and
recomputes $I$. It performs no floating-point root approximation and does not
trust a solver's assertion that all roots were found.

This enumeration is intentionally simple and sometimes inefficient. Large
coefficients can make $B$ enormous. The prototype refuses bounds above
$100\,000$, returning unknown from the producer. That is a resource policy,
not a root-existence conclusion. A later positive-factor or Sturm cover could
reduce work. Forge already proposes Sturm machinery, so reusing a proved
version is preferable to declaring another root-counting implementation new.

\status{The 47 executed singularity objects verify complete seed \emph{plans}
for their encoded operators. They do not contain proofs that two arbitrary
sequences satisfy the recurrences, and they do not prove the seed equalities.}

\subsection{Apparent singularities introduced by transport}

A second recorded example starts from
$A=E-(n+1)$ and $B=E-2$. The producer returns a polynomial common multiple with
leading coefficient $n-1$:
\[
L=(n-1)E^2+(-n^2-3n+2)E+2n(n+1).
\]
Its multipliers are
\[
U=(n-1)E-2n,\qquad V=(n-1)E-n(n+1).
\]
Although both original operators are monic, this common multiple has a natural
singularity at $1$. The safe equality plan includes the additional seed at
$1+2=3$. One may be able to choose another operator or exploit the originals to
avoid that extra work. One may not simply discard the factor or omit the seed
because the singularity ``came from the algorithm''.

This observation gives the scheduler a concrete optimisation criterion:
compare candidate common multiples by checked operator size, estimated
verification cost, and the number and difficulty of their residual seed goals.
It is not necessary to introduce a new global scheduling architecture to use
that cost signal.

\subsection{Connecting all three workers}

For a unary rational matrix $M$, the rows
$u,uM,\ldots,uM^d$ are linearly dependent. Exact search can propose scalars
$c_i$ with $\sum_i c_i uM^i=0$. Multiplication by $M^nv$ gives a
constant-coefficient recurrence for the output sequence. This is another small
identity certificate, and can also be obtained through a characteristic
polynomial. That extraction is a proposed extension, not part of the recorded
runner.

A hypergeometric sum can independently receive a telescoping recurrence.
A claimed closed form, generating-function coefficient sequence, or library
sequence may carry yet another annihilator. Their verified operators meet
through \eqref{eq:left-multiple}; Theorem~\ref{thm:seeds} supplies the final
obligations. In the simplest squared-binomial example, both sides already share
a first-order operator, so common-multiple search should be skipped.

The resulting proof plan is explicit:
\[
\begin{gathered}
\text{source correspondence}\;\longrightarrow\;
\text{two proved annihilation statements}\\
\longrightarrow\;\text{checked common operator}\;
\longrightarrow\;\text{complete singularity cover}\\
\longrightarrow\;\text{proved seed equalities}\;
\longrightarrow\;\forall n,\ f(n)=g(n).
\end{gathered}
\]
Each arrow has an independent obligation. No stage is allowed to stand in for
one of the stages after it.
